% Pursuing Stacks (cotes 134-2 to 134-8): archive cote 134-7, site page 41 (PDF page 42), typescript page 483.
% typescript. Draft compared: ps6.tex:1985-2059.
% AI-transcribed by Claude Opus 5.5 (claude-opus-5-5, adaptive thinking, effort high), Message Batches API, batch ps3 (msgbatch_01XdFmuVVBAcgc7EFxdtCqsj), 2026-11-12; not yet checked by eye.
\scanpage{134-7/41}
\mspage{483}

(4) \[ \mathrm{Hom}_{\underline{O}_k}( \underline{O}_k^{(X)}, F ) \hand{$\simeq$} \mathrm{Hom}(X,F\hand{)} \hand{(}\deleted{\unsure{def}}\hand{$\simeq$} \Gamma(X,F_X)\hand{)}, \]
where in the last member (included as a more geometric interpretation of
the second) $F_X$ denotes the restriction of $F$ to the object $X$, more accu-%
rately to the topos (or site) induced \xout{up}on $X$ b\hand{y} the ambient topos (or
site) we are working in. Thus we may indeed view (1) as the most per-%
fect notion of linearization, as far as generality goes - it makes
sense of course in any ringed topos (without even a commutativity assump-%
tion !). The only trouble is that, even for such a down to earth $X$ as
a unipotent bundle, \xout{say} the standard affine line $E_k^1$ say, the sheaf
$\underline{O}_k^{(X)}$ in (1) is not quasi-coherent and therefore not too amenable it
seems to computations - thus we get easily from (4) a canonical map
(for general $X$)
\[ k^{(X(k))} \hand{$\to$} \Gamma(k, \underline{O}_k^{(X)}) \ (= \underline{O}_k^{(X)}\ (k)) \]
(where the $\Gamma$ in the second member denotes sections over $k$ i.e. value of
a functor on $\mathrm{Alg}_{/k}$ on the initial object $k$ , and remembering in the
first member that the ring of s ections of $\underline{O}_k$ is $k$ ), but I would be at
a loss to make a guess \xout{when}\inserted[typed]{as} for reasonable conditions for this map to
be an isomorphism ! This may seem a prohibitive ``contra'' against using
at all such huge sheaves as $\underline{O}_k^{(X)}$ , the point though is that in most
questions where such linearizations are introduced (mainly questions
where interest lies in computing cohomology invariants), one is practi-%
cally never interested in taking the groups of sections of these, but
rather in looking at their maps into sheaves of modules $F$ precisely,
which is achieved by (4), or taking more generally their $\mathrm{Ext}^i$ with
such an $F$, which is achieved by the similar familiar formula

(5) \[ \mathrm{Ext}^i_{\underline{O}_k}(\underline{O}_k^{(X)}, F ) \hand{$\simeq$} H^i(X,F_X) \quad , \]
more neatly

(5') \[ \mathrm{RHom}_{\underline{O}_k}( \underline{O}_k^{(X)}, F ) \hand{$\simeq$} \mathrm{R}\Gamma_X(F_X) \quad , \]
valid of course again for any ringed topos. In the present context
however, the ``coefficients'' $F$ we are interested in, as was just pointed
out, are not arbitrary $\underline{O}_k$-Modules, but rather \uline{quasi}-\uline{coherent} ones.
Thus, if we get a variant of (1)

(6) \[ X \hand{$\longrightarrow$} L(X) \]
with $L(X)$ some \uline{quasi}-\uline{coherent} Module, giving rise to (4), or even to
(5) and (5'), this would be for us a perfectly good substitue for (3),
which would deserve the name of a ``\uline{quasicoherent} \uline{envelope}'' of $X$. Of

\margin{41}
